\loai{Năng lượng toả ra}
\begin{vidu}%[3P7-7.5K][Loai4]
\nguon{QG - 2017} Rađi $\Hn{88}{226}{Ra}$ là nguyên tố phóng xạ $\alpha $. Một hạt nhân $\Hn{88}{226}{Ra}$ đang đứng yên phóng xạ ra hạt $\alpha $ và biến đổi thành hạt nhân con X. Biết động năng của hạt $\alpha $ là 4,8 MeV. Lấy khối lượng hạt nhân (tính theo đơn vị u) bằng số khối của nó. Giả sử phóng xạ này không kèm theo bức xạ gamma. Năng lượng tỏa ra trong phân rã này là
\choice
{271 MeV}
{4,72 MeV} 
{\True 4,89 MeV} 
{269 MeV}
\loigiai{
\begin{itemize}
\item Phương trình phản ứng: $\Hn{88}{226}{Ra}\rightarrow \Hn{2}{4}{He}+\Hn{86}{222}{Rn}$.
\item Định luật bảo toàn năng lượng toàn phần: $W=K_{Rn}+K_\alpha\quad (1)$
\item Ta có: $\dfrac{K_{Rn}}{K_\alpha}=\dfrac{m_\alpha}{m_{Rn}}=\dfrac{4}{222} \Rightarrow K_{Rn}=\dfrac{4}{222}.K_\alpha=\dfrac{4}{222}.4,8=\dfrac{16}{185}\ MeV$.
\item Thay vào (1) ta có: $W=4,8+\dfrac{16}{185}\approx 4,89\ MeV$.
\end{itemize}
}
\end{vidu}
\loai{Động năng hạt phóng xạ}
\begin{vidu}%[3P7-7.5K][Loai5]
Cho phản ứng phóng xạ: $\Hn{54}{210}{Po}\to \Hn{2}{4}{He}+X+W$ trong đó $X$ là hạt nhân con và $W$ là năng lượng tỏa ra từ phản ứng. Cho biết có thể lấy gần đúng khối lượng của 1 hạt nhân (theo đơn vị $u$) bằng số khối của nó. Giả sử phóng xạ này không kèm theo bức xạ gamma. Động năng của hạt $\alpha$ sinh ra bằng
\choice
{\True 0,98W}
{0,02W}
{W}
{0,5W}
\loigiai{
\begin{itemize}
\item Định luật bảo toàn năng lượng toàn phần: $W=K_{X}+K_\alpha\quad (1)$.
\item Ta có: $\dfrac{K_{X}}{K_\alpha}=\dfrac{m_\alpha}{m_{X}}=\dfrac{4}{206}=\dfrac{W-K_\alpha}{K_\alpha} \Rightarrow 4K_\alpha=206W-206K_\alpha \Rightarrow K_\alpha=\dfrac{206}{210}W\approx 0,98W$.
\end{itemize}	
}
\end{vidu}
\loai{Tốc độ hạt phóng xạ}
\begin{vidu*}%[3P7-7.5G][Loai6]
Hạt nhân Poloni đứng yên, phóng xạ $\alpha$ biến thành hạt nhân X. Cho $m_{Po}=209,9373u;\ m_{\alpha}=4,0015u;\ {m_X}=205,9294u$. Giả sử phóng xạ này không kèm theo bức xạ gamma. Biết $1u=1,66055.10^{-27}\ kg$. Tốc độ hạt $\alpha$ phóng ra là
\choice
{$1,27.10^{7}\ m/s$}
{\True $1,68.10^{7}\ m/s$}
{$2,12.10^{7}\ m/s$}
{$3,27.10^{7}\ m/s$}
\loigiai{
\begin{itemize}
\item $p_{\alpha}={p_X}\Rightarrow m_{\alpha}v_{\alpha}={m_X}{v_X}\Rightarrow 4,0015v_{\alpha}=205,9294{v_X}\Rightarrow {v_X}=0,01943v_{\alpha}\quad (1)$.
\item Năng lượng tỏa ra của phản ứng: $W=(m_{Po}-m_{\alpha}-{m_X})c^2=5,9616\ MeV=9,539.10^{-13}\ J$.
\item Áp dụng bảo toàn năng lượng ta có:
\begin{align*}
&W=K_{\alpha}+K_X=\dfrac{1}{2}m_{\alpha}v_{\alpha}^2+\dfrac{1}{2}{m_X}v_{X}^2.\\
\Rightarrow\ &W=9,539.10^{-13}=\dfrac{1}{2}. 4,0015. 1,66055.10^{-27}v_{\alpha}^2+\dfrac{1}{2}. 205,9294. 1,66055.10^{-27}v_{X}^2\quad (2).
\end{align*}
\item Từ $(1)$ và $(2)$ suy ra $v_{\alpha}=1,68.10^{7}\ m/s$.
\end{itemize}
}
\end{vidu*}
\loai{Động năng của hạt nhân con}
\begin{vidu}%[3P7-7.5K][Loai7]
Hạt nhân $^{234}U$ đứng yên phân rã $\alpha$ biến đổi thành hạt nhân $X$. Biết khối lượng của các hạt nhân: $m_U=233,9905u$, $m_\alpha =4,0015u$, $m_X=229,9838u$. Giả sử phóng xạ này không kèm theo bức xạ gamma. Lấy $1u=931,5\ MeV/c^2$. Hạt nhân $X$ giật lùi với động năng bằng
\choice
{$82,8 KeV$}
{$4,76\ MeV$}
{\True $0,0828\ MeV$}
{$47,6\ KMeV$}
\loigiai{
\begin{itemize}
\item Tổng động năng của hai hạt sau phản ứng: $K_X + K_\alpha =W=4,84\ MeV\quad (1)$.
\item Động lượng của hệ bảo toàn nên ta có:
$p_\alpha =p_X \Rightarrow 2m_\alpha K_\alpha = 2 m_XK_X\Rightarrow K_\alpha = 57,5K_X\quad (2)$.
\item Từ $(1)$ và $(2)$ ta tìm được: $K_X=0,0828\ MeV$.
\end{itemize}
}
\end{vidu}
\loai{Tốc độ hạt nhân con}
\begin{vidu}%[3P7-7.5K][Loai8]
\nguon{ĐH-2012} Một hạt nhân X, ban đầu đứng yên, phóng xạ $\alpha$ và biến thành hạt nhân Y. Biết hạt nhân X có số khối là A, hạt $\alpha$ phát ra tốc độ v. Lấy khối lượng của hạt nhân bằng số khối của nó tính theo đơn vị u. Tốc độ của hạt nhân Y bằng
\choice
{$\dfrac{4v}{A+4}$}
{$\dfrac{2v}{A-4}$}
{\True $\dfrac{4v}{A-4}$}
{$\dfrac{2v}{A+4}$}
\loigiai{
Áp dụng định luật bảo toàn động lượng $(A - 4)\vec{V}+ 4\vec{v} = 0 \Rightarrow \vec{V}=\dfrac{-4\vec{v}}{A-4}\Rightarrow $ Độ lớn $V = \dfrac{4v}{A-4}$.
}
\end{vidu}
\loai{Bài toán tổng hợp}
\begin{vidu}%[3P7-7.5K][Loai9]
Hạt nhân $\Hn{84}{210}{Po}$ đứng yên, phân rã $\alpha$ thành hạt nhân chì. Giả sử phóng xạ này không kèm theo bức xạ gamma.  Lấy khối lượng nguyên tử tính theo đơn vị u bằng số khối của hạt nhân của nguyên tử đó. Động năng của hạt $\alpha$ bay ra bằng bao nhiêu phần trăm của năng lượng phân rã?
\choice
{\True 98,1\%}
{13,8\%}
{1,9\%}
{86,2\%}
\loigiai{
\begin{itemize}
\item Phương trình phản ứng: $\Hn{84}{210}{Po}\to \alpha +\Hn{82}{206}{Pb}$.
\item Ta có: $K_{Pb}=\dfrac{m_\alpha}{m_{Pb}}K_\alpha $.
\item Theo định luật bảo toàn năng lượng ta có:
\begin{align*}
W=K_\alpha +K_{Pb}=K_\alpha +\dfrac{m_\alpha}{m_{Pb}}K_\alpha =\dfrac{m_{Pb}+m_\alpha}{m_{Pb}}K_\alpha \Leftrightarrow \dfrac{K_\alpha}{W}=\dfrac{m_{Pb}}{m_{Pb}+m_\alpha}=\dfrac{206}{206+4}=0,981\left({98,1\%}\right).
\end{align*}
\end{itemize}
}
\end{vidu}
\begin{vidu}%[3P7-7.5K][Loai9]
Hạt nhân mẹ $Ra$ đứng yên biến đổi thành một hạt $\alpha$ và một hạt nhân con $Rn$. Giả sử phóng xạ này không kèm theo bức xạ gamma. Tính động năng của hạt $\alpha$ và hạt nhân $Rn$. Biết $m_{Ra}=225,977u$, $m_{Rn}=221,970u$; $m_\alpha=4,0015u$.
\choice
{$K_\alpha =0,09\ MeV$; $K_{Rn}=5,03\ MeV$}
{$K_\alpha =0,009\ MeV$; $K_{Rn}=5,3\ MeV$}
{$K_\alpha =503\ MeV$; $K_{Rn}=90\ MeV$}
{\True $K_\alpha =5,03\ MeV$; $K_{Rn}=0,09\ MeV$}
\loigiai{
\begin{itemize}
\item Phương trình phản ứng hạt nhân là $Ra\to \alpha +Rn$.
\item Năng lượng của phản ứng hạt nhân tỏa ra là:
\begin{align*}
\left\{\begin{aligned}& W=\left({m_{Ra}-m_\alpha-m_{Rn}}\right)c^2=K_{Rn}+K_\alpha=5,123\ MeV\\
& p_\alpha =p_{Rn}\Leftrightarrow m_\alpha K_\alpha = m_{Rn}K_{Rn}
\end{aligned}\right. \Rightarrow \left\{\begin{aligned}& K_{Rn}=0,09\ MeV. \\
& K_\alpha =5,03\ MeV. 
\end{aligned}\right.
\end{align*}
\end{itemize}
}
\end{vidu}
